\subsection{A clock read from source interval counts}
\label{sec:source-count-clock}

The equal-mass source populations of
Proposition~\ref{prop:golden-source-count-limit} admit a clock defined by
their records. Suppose an observer can read the complete ancestry of an
upper event and identify a lower anchor. Its inclusive interval is
$I(p,q)=\{e:p\preceq e\preceq q\}$. Choose a nondegenerate reference
interval $J$ in the same history and define
\begin{equation}
 \mathcal T(I;J)=\left(\frac{|I|}{|J|}\right)^{1/4}.
 \label{eq:source-count-clock}
\end{equation}
This readout uses order and cardinality. It needs no numerical event
labels, spatial coordinates, timestamps, absolute density or oscillator
frequency. Complete ancestry access is a substantive readout hypothesis:
a hash of an unavailable record does not provide its contents. The cost
of collecting and counting the records is separate from the duration
reconstructed by the readout.

\begin{theorem}[Count clocks and proper-time ratios]
\label{thm:source-count-clock}
Consider the declared source histories with their controlled order and
equal-mass count limit. Let $I_n,J_n$ approximate fixed interior timelike
diamonds, with one common event weight $w_n>0$, and let their limiting
proper-time separations be $\tau_I,\tau_J>0$. Then
\[
 \mathcal T(I_n;J_n)\longrightarrow\frac{\tau_I}{\tau_J}.
\]
The reference interval fixes the time unit; the clock origin is free.
\end{theorem}
\begin{proof}
For coordinate measure $dt\,d^3x$ and metric $c^2dt^2-|dx|^2$, a flat
diamond of proper-time separation $\tau$ has volume
$K\tau^4$, where $K=\pi c^3/24$. The vertical formula follows by
integrating its ball sections. A Lorentz transformation has determinant
one and sends a timelike tip displacement to its rest form, so the same
formula holds for moving tips. The count-volume theorem gives
$w_n|I_n|\to K\tau_I^4$ and $w_n|J_n|\to K\tau_J^4$.
Taking their ratio and its positive fourth root proves the claim. Both
the event weight and the volume constant cancel.
\end{proof}

The supplied relation $\Delta=a/c$ enters the proof of the order and
volume limit. It is absent from the clock's definition. This distinction
allows a retrospective clock to be read from the same history without
assuming an oscillator Hamiltonian. It does not select the population,
read law or laboratory unit from accepted microscopic repairs.

\begin{proposition}[Finite clock enclosure]
\label{prop:source-count-clock-error}
Suppose $A=wN$, $B=wM$ and independently justified volume bounds obey
$|A-V_I|\leq E_I$, $|B-V_J|\leq E_J$, with $B>E_J\geq0$,
$E_I\geq0$ and $V_J>0$. Then
\begin{equation}
 \left(\frac{(A-E_I)_+}{B+E_J}\right)^{1/4}
 \leq\frac{\tau_I}{\tau_J}\leq
 \left(\frac{A+E_I}{B-E_J}\right)^{1/4}.
 \label{eq:source-count-clock-error}
\end{equation}
The same formula uses counts $N,M$ and errors $E_I/w,E_J/w$.
\end{proposition}
\begin{proof}
The true volumes belong to their stated intervals. Division by a positive
denominator and the increasing fourth root preserve the two inequalities.
The error bound in \eqref{eq:source-net-general-volume-error} supplies
such volume errors when its geometric hypotheses hold. Counts alone do
not supply those errors.
\end{proof}

For relative volume errors $A=V_I(1+d_I)$ and $B=V_J(1+d_J)$, with
$|d_I|\leq\eta_I<1$ and $|d_J|\leq\eta_J<1$, the measured clock divided
by the exact proper-time ratio lies in
\begin{equation}
 \left[
 \left(\frac{1-\eta_I}{1+\eta_J}\right)^{1/4},
 \left(\frac{1+\eta_I}{1-\eta_J}\right)^{1/4}
 \right].
 \label{eq:source-count-clock-relative}
\end{equation}

\begin{proposition}[Unit freedom within the volume clock class]
\label{prop:source-count-clock-uniqueness}
A positive elapsed-time function $F(V)$ depending only on diamond volume
and additive for successive collinear inertial intervals is
$F(V)=\kappa V^{1/4}$ for some $\kappa>0$.
\end{proposition}
\begin{proof}
Set $g(t)=F(Kt^4)$ for $t>0$. Inertial additivity gives
$g(t+s)=g(t)+g(s)$. Positivity implies strict monotonicity. Additivity
first fixes $g$ on positive rationals; monotonicity and rational
approximation give $g(t)=g(1)t$ for all positive real $t$.
This is uniqueness within the stated class, not a premise about every
physical pointer.
\end{proof}

\paragraph{Accumulating duration along a curve.}
Let $\gamma(t)=(t,x(t))$ be $C^2$ on $[0,T]$, with
$|x'(t)|\leq v<c$, $|x''(t)|\leq A_0$, and a common spatial buffer for
its small diamonds. Segment durations are comparable to
$\delta_n\to0$. Suppose
\[
 \frac{\Delta_n}{\delta_n}\to0,\qquad
 \frac{H_n}{\delta_n}\to0,\qquad
 \frac{h_n}{a_n}\to0,
\]
and sampled tips have spatial and temporal errors $O(H_n+\Delta_n)$.
Summing \eqref{eq:source-count-clock} over successive segment intervals,
using the same fixed reference interval, converges to
\begin{equation}
 \frac1{\tau_J}\int_0^T\sqrt{1-|x'(t)|^2/c^2}\,dt.
 \label{eq:source-count-clock-curve}
\end{equation}
Indeed, the two-shell estimate is uniformly
$O(\delta_n^4h_n/a_n+\delta_n^3(H_n+\Delta_n))$ per segment.
The strict speed margin bounds each exact diamond volume below by a
positive constant times $\delta_n^4$. Tip perturbations have the same
relative error order. Equation~\eqref{eq:source-count-clock-relative}
therefore controls all segment clocks by one vanishing relative error.

The continuum polygon error has an explicit bound. For
$f(u)=\sqrt{1-|u|^2/c^2}$ on $|u|\leq v$, its Hessian norm is at most
$M_0=[c^2(1-v^2/c^2)^{3/2}]^{-1}$. Taylor expansion about each segment's
average velocity has zero integrated linear term. Since the velocity
differs from that average by at most $A_0\delta_i$, the chord and curve
proper times differ by at most $M_0A_0^2\delta_i^3/2$. The total is at
most $M_0A_0^2T(\max_i\delta_i)^2/2$. This proves
\eqref{eq:source-count-clock-curve}. Fixed one-layer intervals do not
satisfy these resolution hypotheses.

\paragraph{Finite records and calibration scope.}
The finite executable control parses writer identities, immutable
register versions and values in the $q=5$ history. Its $500$ events give
the three centered inclusive counts $2,41,80$. The first two continuum
diamonds lie inside the source window; the third does not. Exact dyadic
fourth-root bounds, schedule permutation and opaque event relabeling
check the readout. These coarse intervals provide no finite proper-time
accuracy certificate. The source-pinned expected hash authenticates the
replayed content relative to that commitment; it is not an external
signature or a proof of physical provenance. Unrelated serial audit
links are excluded from the semantic read order.

Lean formalizes the finite interval and root inequalities. The count-volume
and curve limits above are analytic proofs. In this layered law every causal pair $k$
layers apart has a path of exactly $k$ reads, and every strict chain has
at most $k$ edges. Longest-chain height consequently reads coordinate
duration independently of displacement. It is not a proper-time clock
for moving tips in this family. The count construction concerns
geometric duration reconstructed from accessible records; it neither
sets the time cost of primitive seam words nor identifies physical clock
readings without a common physical interpretation of the event law.
